\subsection{Dense original good opens over regular centers}
\begin{lemma}[Principal opens of regular elements are dense]
\label{scheme:regular-principal-open-dense}
\lean{EnvelopeProperty.dense_basicOpen_of_mem_nonZeroDivisors}
\leanok
Let $R$ be any commutative ring and let $r\in R$ be a non-zero-divisor.
Then the actual principal open $D(r)$ is dense in $\operatorname{Spec}R$.
\end{lemma}
\begin{proof}
Every power of $r$ is a non-zero-divisor, so the original localization
map $R\to R[1/r]$ is injective. The closure of its actual spectrum image
is the zero locus of its kernel, hence the whole spectrum. That image
is exactly $D(r)$. This argument also includes the zero ring.
\end{proof}
\begin{theorem}[Both original Rees good opens are dense]
\label{scheme:actual-rees-good-open-density}
\lean{EnvelopeProperty.ReesPresentation.coefficient_mem_nonZeroDivisors,EnvelopeProperty.ReesPresentation.dense_projection_preimage_basicOpen}
\leanok
\uses{scheme:actual-rees-good-open-isomorphism,scheme:regular-principal-open-dense,scheme:actual-rees-coordinate-cover}
Let $A$ be any commutative ring, let $(g_i)_{i\in I}$ be any family,
and suppose the original $g_j$ is a non-zero-divisor in $A$.
Then $D(g_j)$ is dense in the original target and its full inverse image
is dense in the SAME original Rees scheme. The actual projection is
therefore an isomorphism between these genuine dense opens.
\end{theorem}
\begin{proof}
On EVERY original coordinate chart indexed by $i$, the genuine whole
chart ring injects into the original $A[1/g_i]$. Localization preserves
regularity of the original $g_j$, and injectivity reflects it back to
its original coefficient image in the chart. Thus the preceding lemma
makes the actual principal open of that coefficient dense in every
original chart. The whole coefficient projection triangle identifies
this open with the inverse image of the original target $D(g_j)$.
For any nonempty open subset of the Rees scheme, choose an actual point
and a member of the genuine coordinate affine open cover above that point.
Its open inverse image is nonempty, so it meets the dense chart principal
open. Mapping that point back gives a point of the desired intersection
in the SAME original scheme. Target density follows from the lemma,
and the previously proved actual good-open pullback isomorphism gives
the asserted isomorphism. No integrality, reducedness, Noetherianity,
finite generation, or supplied birationality comparison is used here.
\end{proof}
\begin{theorem}[Surjectivity for a regular original center]
\label{scheme:actual-rees-regular-center-surjective}
\lean{EnvelopeProperty.ReesPresentation.basicOpen_subset_range_projection,EnvelopeProperty.ReesPresentation.projection_surjective_of_mem_nonZeroDivisors}
\leanok
\uses{scheme:actual-rees-good-open-density,scheme:actual-rees-presentation-proper}
If the original family is finite and one of its original generators is
a non-zero-divisor, the SAME original proper Rees projection is surjective.
\end{theorem}
\begin{proof}
The actual good-open lift shows that the image contains all of the
original $D(g_j)$. That open is dense by the preceding result. The proved
properness of the SAME projection makes its actual whole image closed.
A closed set containing this dense open is the entire original target.
Thus every original target point has a genuine source point. No domain
assumption or point-lifting witness is supplied.
\end{proof}
